\minitoc 

% ===================================================================================================================
% Introduction 


% ===================================================================================================================
% Cadre Statistique

\section{Cadre}

Dans cette section, nous posons le cadre de l'étude des statistiques ainsi que l'ensemble 
des notations et objets utilisés. 

\subsection{Modèle}


L'étude des statistiques se résume à l'étude de la loi de variables aléatoires ou de suites de 
variables aléatoires. Nous devons donc définir un cadre dans lequel ces différents objets vivrons. 
On considère ainsi des variables aléatoires $X$ définies telles que : 
\begin{equation}
    X : 
        \begin{cases}
            ( \Omega, \mathcal{T}, \myP) \longrightarrow (\mathcal{X}, \mathcal{A}, \myP^X) \\ 
            \omega \longmapsto X( \omega)
        \end{cases}    
\end{equation}

Dans notre contexte, nous n'avons accès qu'à des \emph{réalisations} $X( \omega)$ de la variable 
aléatoire $X$ et nous souhaitons déterminer sa loi $ \myP^X$. 
Pour cela, nous faisons une hypothèse appelée \emph{hypothèse paramétrique forte} qui suppose que 
la loi de $X$ ne dépend que d'un paramètre $ \theta \in \Theta$, où $ \Theta \subset \R^d$. 
Ainsi, la recherche de la loi $ \myP^X$ se réduit à la détermination d'un paramètre $ \theta_*$ tel 
que $ \myP^X = \myP_{\theta_*}$. 

\begin{definition}[Modèle Paramétrique]
    On appellera \emph{modèle paramétrique} l'espace probabilisé $( \mathcal{X}, \mathcal{A}, \mathcal{P}_\Theta)$ 
    tel que $ \mathcal{P}_\Theta = \{ \myP_\theta \mid \theta \in \Theta\}$ où $ \Theta$ est un ouvert convexe 
    de $ \R^d$. 
\end{definition}

Or cet espace ne nous permet de considérer que de simples variables aléatoires. Dans un cadre plus général, 
nous souhaiterions étudier des \emph{échantillons} $(X_1, \dots, X_n)$ i.i.d d'une variable aléatoire 
$X$. Nous définissons pour cela le \emph{modèle d'échantillonnage}. 

\begin{definition}[Modèle d'échantillonnage]
    Soit $( \mathcal{X}, \mathcal{A}, \mathcal{P}_{\Theta})$ un modèle paramétrique. 
    Un \emph{échantillon} $(X_1( \omega), \dots, X_n( \omega))$ de ce modèle est un $n$-uplet de réalisations 
    indépendantes d'une variable aléatoire $X$ de loi $ \myP_\theta \in \mathcal{P}_\Theta$. 
    On note $ \myP^{\{X_1, \dots, X_n\}} = \myP_{\theta_*}^{ \otimes n}$ la loi de cet échantillon. Le modèle 
    paramétrique de cet échantillon est donc : 
        $$ ( \mathcal{X}, \mathcal{A}, \mathcal{P}_\Theta)^{\otimes n} = ( \mathcal{X}^n, \mathcal{A}^n, \mathcal{P}_{\Theta}^{\otimes n}) $$
\end{definition}

On a donc : 
    $$ 
        (X_1, \dots, X_n) : 
            \begin{cases}
                ( \Omega, \mathcal{T}, \myP) &\longrightarrow ( \mathcal{X}, \mathcal{A}, \mathcal{P}_\Theta)^{\otimes n} \\ 
                \omega &\longmapsto (X_1( \omega), \dots, X_n( \omega))
            \end{cases}
    $$

Enfin, pour l'étude plus théorique de tels objets, il serait intéressant de pouvoir faire 
tendre vers l'infini la taille de ces échantillons $(X_1, \dots, X_n)$. On définit donc un 
modèle dit \emph{asymptotique} : 
    $$ 
        (X_i)_{i \in \N^*} : 
        \begin{cases}
            ( \Omega, \mathcal{T}, \myP) & \longrightarrow ( \mathcal{X}, \mathcal{A}, \mathcal{P}_\Theta)^\infty \\ 
            \omega & \longmapsto \{X_i( \omega) \mid i \in \N_*\}
        \end{cases}
    $$


\subsection{Les variables aléatoires d'un modèle}

L'objet principal que nous allons manipuler permet d'effectuer des opérations sur des échantillons 
d'une variable aléatoire. On les appelle les \emph{statistiques}, notées $S$. 

\begin{definition}[Statistique]
    Une fonction mesurable $S$ est une \emph{statistique} si elle ne dépend pas 
    de $ \theta$. On a donc : 
        $$ 
            S(X_1, \dots, X_n) : 
                \begin{cases}
                    ( \mathcal{X}, \mathcal{A})^n \longrightarrow ( \mathcal{Y}, \mathcal{B}) \\ 
                    (X_1, \dots, X_n) \longmapsto S(X_1, \dots, X_n)
                \end{cases}
        $$ 
    On notera aussi $ \myP_\theta^S = ( \myP_\theta^{\oplus n})^S$ la mesure image de 
    $\myP_\theta^{\oplus n}$ par $S$. 
\end{definition}

Cette translation de nos observations $(X_1, \dots, X_n)$ implique la définition d'un \emph{modèle image} 
nous permettant de manipuler ce nouvel objet $S(X_1, \dots, X_n)$. 

\begin{definition}[Modèle Image]
    Le modèle image de $( \mathcal{X}, \mathcal{A}, \mathcal{P}_\Theta)^{\oplus n}$ par $S_n$ 
    est donc $ ( \mathcal{Y}, \mathcal{B}, \mathcal{P}_\Theta^{S_n})$, avec : 
        \begin{equation}
            \mathcal{P}_\Theta^{S_n} := \{ \myP_\theta \circ S_n^{-1} \mid \theta \in \Theta\}
        \end{equation}
\end{definition}

Pour rappel, l'objectif des statistiques est de déterminer le paramètre $ \theta_* \in \Theta$. 
Pour cela, il faut que la donnée utilisée ($(X_1, \dots, X_n)$, $X$, ou $S(X_1, \dots, X_n)$) contienne 
le plus d'information possible sur $ \theta$. 
Il serait donc préférable que la translation $S$ ne réduise pas cette quantité. De telles statistiques seront 
donc dites \emph{exhaustives}. Elles seront caractérisées par la loi conditionnelle de l'échantillon. 

\begin{definition}[Loi conditionnelle d'une statistique]
    Soit $T$ une statistique. La loi conditionnelle de $(X_1, \dots, X_n)$ sachant 
    $T(X_1, \dots, X_n)$ est définie par : 
        \begin{equation}
            \myP_\theta^{(X_1, \dots, X_n) \mid T}(A) = \myP((X_1, \dots, X_n) \in A \mid T) 
            = \mathbb{E}_\theta( \mathbbm{1}_{(X_1, \dots, X_n) \in A} \mid T )
        \end{equation}
    pour tout $A \in \mathcal{A}^n$. 
\end{definition}

% \begin{definition}[Statistique Exhaustives]
%     Soit $T$ une statistique sur $( \mathcal{X}, \mathcal{A}, \mathcal{P}_\Theta)^n$ est dite \emph{exhaustive en $ \theta$}
%     si $ \myP_\theta^{(X_1, \dots, X_n) \mid T}$ ne dépend pas de $ \theta$. Cela revient 
%     à dire que l'application : 
%         $$ \theta \longmapsto \frac{ \myP_\theta^n (A \cap T^{-1}(C))}{ \myP_\theta^T(C)} 
%             = \frac{ \myP_\theta((X_1, \dots, X_n) \in A \cap T(X_1, \dots, X_n) \in C)}{ \myP_\theta T(X_1, \dots, X_n) \in C} $$
%     est constante pour tout $ \theta \in \{ \theta \in \Theta \mid \myP_\theta^T(C) > 0\}$. 
% \end{definition}

% Une statistique exhaustive est donc une variable aléatoire telle que la loi 
% conditionnelle d'un échantillon par rapport à cette statistique ne dépend pas de $ \theta$. 
% Elles apportent toute l'information nécessaire sur $ \theta^*$ contenue 
% dans l'échantillon. Ce sont donc de bonnes statistiques à utiliser. 

\subsection{Domination et Densités}

\begin{definition}[Domination]
    On dira qu'une mesure $ \mu$ positive \emph{domine} $ \myP_\theta$ sur $( \mathcal{X}, \mathcal{A})$ si : 
    \begin{equation}
        \boxed{
            \forall A \in \mathcal{A}, \quad \mu(A) = 0 \Longrightarrow \myP_\theta(A) = 0 
        }
    \end{equation}
    On notera alors $ \myP_\theta \ll \mu$.
    Si c'est le cas pour tout $ \theta \in \Theta$, on dira alors que \emph{le modèle est dominé par $ \mu$}.
\end{definition}

Plus formellement, $ \myP_\theta \ll \mu$ ssi 
    \begin{equation}
        \forall \varepsilon > 0, \exists \delta > 0 \quad \text{tel que} \quad 
        \forall A \in \mathcal{A}, \mu(A) < \delta \Longrightarrow \myP_\theta(A) < \varepsilon 
    \end{equation}

\begin{theorem}[Radon-Nikodym]
    Soit $ \mu$ $ \sigma$-finie. On a $ \myP_\theta \ll \mu$ ssi 
    \begin{equation}
        \forall \theta \in \Theta, \exists f \in L^1_+( \mu) \quad \text{telle que} \quad 
        \forall A \in \mathcal{A}, \quad \myP_\theta(A) = \int_A f \; d \mu 
    \end{equation}
    La densite $f$ de $ \myP_\theta$ est unique à égalité $ \mu$-presque partout. Elle est notée 
    $f = d \myP_\theta / d \mu$ comme dérivée de Radon-Nicodym. 
\end{theorem}

\begin{remark}[Cas discret]
    Le cas des lois discrètes peut paraître contre-intuitif, mais s'explique bien 
    grâce à ce théorème. Soit $ \mathcal{X}$ fini ou dénombrable et $ \mathcal{A} = \mathcal{P}( \mathcal{X})$. 
    Soit $ m_c$ la mesure de comptage sur $ \mathcal{A}$. Par définition de l'intégrale de Lebesgue, pour toute 
    fonction $ f : \mathcal{X} \longrightarrow \R$ et $ A \in \mathcal{A}$, on a : 
        $$ \int_A f d \; m_c = \sum_{x \in A} f(x) m_c(x) = \sum_{x \in A} f(x) $$ 
    Soit $ \mathcal{P}_\Theta$ une famille de lois sur $ \mathcal{X}$ et $n = 1$. On a $ \mathcal{P}_\theta \ll m_c$ et 
    les densités $f = d \myP_\theta / d m_c$ existent et vérifient :
        $$ \forall A = {x} \in \mathcal{A}, \quad \myP_\theta(\{x\}) = \int_{\{x\}} f \; d m_c = f(x) $$
    On a donc bien $f(x) = \myP_\theta(X = x)$. 
\end{remark}

\subsection{Hypothèses de régularité}

Pour tous les raisonnements suivants sur notre modèle, nous devons définir des \emph{hypothèses}. 
Celles-ci nous assureront certaines propriétés nécessaires à l'application des théorèmes suivants. 


\begin{definition}[Modèle Identifiable]
    S'il existe $ \mu$ $ \sigma$-finie telle que $ \mathcal{P}_\Theta \ll \mu$ et $ \mu(\{f_{\theta_0} \not = f_{\theta_1}\}) > 0$ 
    pour tout $ \theta_0 \not = \theta_1$, alors le modèle $( \mathcal{X}, \mathcal{A}, \mathcal{P}_\Theta)$ 
    est dit \emph{identifiable}. 
\end{definition}

Un modèle est identifiable si $ \theta$ caractérise de façon unique ue loi. Dans le cas contraire, 
un même valeur de $ \theta$ pourrait correspondre à plusieurs lois initiales. L'estimation 
du paramètre n'aurait donc plus aucun sens. 

\begin{definition}[Modèle Homogène]
    Le modèle $( \mathcal{X}, \mathcal{A}, \mathcal{P}_\Theta)$ est dit \emph{homogène} si 
    $ \mathcal{P}_\Theta \ll \myP_\theta$ pour tout $ \theta \in \Theta$. 
\end{definition}

L'homogénéité est une propriété fondamentale d'un modèle et elle intervient sur 
plusieurs aspects. Un modèle homogène garantit que tous les $ \myP_\theta$ ont les mêmes 
ensembles négligeables. D'autre part, l'homogénéité garantit que le support des variables 
aléatoires ne dépend pas de $ \theta$. Sans cela, les inversion de dérivées/intégrales 
(voir ci-dessous) ne seraient pas possibles. 

\begin{proposition}
    Le modèle $( \mathcal{X}, \mathcal{A}, \mathcal{P}_\Theta)$ est homogène ssi, il existe $ \mu$ 
    $ \sigma$-finie telle que : 
    \begin{enumerate}[label=\roman*)]
        \item $ \mathcal{P}_\Theta \ll \mu$ 
        \item $\{f_\theta > 0\}$ $ \mu$-presque partout pour tout $ \theta \in \Theta$. 
    \end{enumerate}
\end{proposition}

En conséquence, pour tout $ \theta \in \Theta$, $ \mathcal{P}_\Theta \ll \myP_\theta$. Donc 
les $ \myP_\theta$ ont tous les mêmes négligeables et la densité de $ \myP_{ \theta_2}$ par rapport 
à $ \myP_{ \theta_1}$ s'écrit : 
    $$ 
        \frac{ d \myP_{ \theta_2}}{ d \myP_{ \theta_1}} = \frac{ f_{ \theta_2}}{f_{ \theta_1}} 
        = \frac{ d \myP_{ \theta_2} / d \mu}{ d \myP_{ \theta_1} / d \mu}
    $$ 
En effet, on a pour tout $A \in \mathcal{A}$ : 
    \begin{align*}
        \myP_{ \theta_2}(A) &= \int_A \frac{f_{ \theta_2}}{f_{ \theta_1}} \; d \myP_{ \theta_1} 
        = \int_A \frac{f_{ \theta_2}}{f_{ \theta_1}} f_{ \theta_1} \; d \mu \\ 
        &= \int_A f_{ \theta_2} \; d \mu = \myP_{\theta_2}(A). 
    \end{align*}

\begin{definition}[Hypothèses de régularité]
    Définissons donc ces hypothèses : 
    \begin{itemize}
        \item[$(H_1)$] : Le modèle est \textbf{identifiable} et \textbf{dominé}. 
        $$ i.e \quad \mathcal{P}_\Theta << \mu, \quad \text{et} \quad \forall \theta \in \Theta, f_\theta = \frac{ d \myP_\theta}{ d \mu}, 
        \quad \text{et} \quad \forall \theta_1 \not = \theta_2, \; \mu(\{f_{\theta_1} = f_{\theta_2}\}) = 0 
        $$
        \item[$(H_2)$] : Le modèle est \textbf{homogène}. 
        $$ 
            i.e \quad \forall \theta \in \Theta, \; \nabla f_\theta \in L^2( \mu)
        $$ 
        \item[$(H_3)$] : Pour tout $A \in \mathcal{A}$, la fonction $ \theta \in \Theta \longmapsto \myP_\theta(A)$, admet des 
            dérivées partielles telles que : 
            $$ 
                \frac{ \partial}{ \partial \theta_k} \myP_\theta(A) = \frac{ \partial}{ \partial \theta_k} \int_A f_\theta \; d \mu
                = \int_A \left( \frac{ \partial}{ \partial \theta_k} f_\theta \right) \;  d \mu 
            $$ 
        \item[$(H_4)$] : Pour tout $A \in \mathcal{A}$, la fonction $ \theta \in \Theta \longmapsto \myP_\theta(A)$, admet des 
            dérivées partielles d'ordre 2 telles que : 
            $$ 
                \frac{ \partial^2}{ \partial \theta_j \partial \theta_k} \myP_\theta(A) = \frac{ \partial^2}{ \partial \theta_j \partial \theta_k} \int_A f_\theta \; d \mu
                = \int_A \left( \frac{ \partial^2}{ \partial \theta_j \partial \theta_k} f_\theta \right) \;  d \mu 
            $$ 
    \end{itemize}
\end{definition}


\subsection{Vraissemblance}

Soit un échantillon ${X} = (X_1, \dots, X_n)$ d'un modèle vérifiant $H_{0,1}$.

\begin{definition}[Fonction de vraissemblance]
    On appelle \emph{fonction de vraissemblance} l'application $L$ définie par : 
        $$ L(\theta; x) = \prod_{i=1}^{n} f_\theta(x_i) $$
    Pour un $x$ observé, $L$ est vue comme une fonction du paramètre $\theta$. Il est crucial d'inclure les 
    indicatrices de support $ \mathbbm{1}_{A(\theta)}(x_i)$ dans l'expression de $f_\theta$.
\end{definition}

\begin{remark}[Log-vraisemblance]
    En pratique, on utilise la log-vraisemblance notée $\ell_n(\theta)$ :
        $$ \ell_n(\theta) = \log L(\theta; x) = \sum_{i=1}^n \log f_\theta(x_i). $$
    Elle facilite la recherche du Maximum de Vraisemblance (EMV) et le calcul du Score.
\end{remark}

\begin{theorem}[Critère de factorisation de Fisher-Neyman]
    $T$ est une statistique exhaustive ssi il existe deux fonctions mesurables positives $h$ et $g_\theta$ telles que :
        $$ \boxed{L(\theta; {x}) = h({x}) \cdot g_\theta(T({x}))} $$
    Où $h$ est indépendante de $\theta$ et $g_\theta$ ne dépend de ${x}$ qu'à travers $T({x})$.
\end{theorem}

% ===================================================================================================================
% Informations Quantitatives

\section{Informations Quantitatives}

\subsection{Exhaustivité}

Une statistique $T$ est exhaustive si elle "conserve" toute l'information utile sur $\theta$. Mathématiquement, 
cela signifie qu'une fois $T$ connue, la loi de l'échantillon ne dépend plus de $\theta$.

\begin{definition}[Statistique Exhaustive]
    $T$ est exhaustive pour $\theta$ si la loi conditionnelle de $(X_1, \dots, X_n)$ sachant $T$ est indépendante de $\theta$. 
    $$ \forall \theta \in \Theta, \quad \mathcal{L}_\theta(X_1, \dots, X_n \mid T) = \mathcal{L}(X_1, \dots, X_n \mid T) $$
\end{definition}

\begin{remark}
    En pratique, on utilise presque exclusivement le {théorème de factorisation de Fisher-Neyman} pour 
    prouver l'exhaustivité.
\end{remark}

\begin{proposition}
    $T$ est une \emph{statistique exhaustive} ssi pour toute statistique $S$ réelle intégrable, on a : 
        \begin{equation}
            \mathbb{E}_\theta(S \mid \sigma(T)) = \varphi( T)
        \end{equation}
    (i.e si cela ne dépend pas de $ \theta$.)
\end{proposition}

\begin{example}[Statistique exhaustive]
    Soit $(X_1, \dots, X_n)$ un échantillon d'une variable aléatoire $ X \sim \mathcal{B}(1, \theta)$. Posons 
    $S_n = \sum_{i=1}^{n} X_i \sim \mathcal{B}(n, \theta)$. On a ici, $ \mathcal{X} = \{0, 1\}^n$. 
    On cherche à estimer $ \theta$. Montrons que $S$ est une statistique exhaustive. 
    Soit $k \in \llbracket 1, n \rrbracket$, on : 
    \begin{align*}
        & \myP_\theta ((X_1, \dots, X_n) = (x_1, \dots, x_n) \mid S_n = k) \\ 
        &= \frac{ \myP_\theta( (X_1, \dots, X_n) = (x_1, \dots, x_n) \cap S_n = k)}{ \myP_\theta(S_n = k)} \\ 
        &= \mathbbm{1}_{\left\{\sum_{i=1}^{n} x_i = k\right\}} \times  \frac{ \myP_\theta( (X_1, \dots, X_n) = (x_1, \dots, x_n))}{ \myP_\theta(S_n = k)} \\ 
        &= \mathbbm{1}_{\left\{\sum_{i=1}^{n} x_i = k\right\}} \times \prod_{i=1}^{n} \frac{ \myP_\theta(X_i = x_i)}{ \binom{n}{k} \theta^k (1 - \theta)^{n - k}} \\ 
        &= \frac{1}{\binom{n}{k}}
    \end{align*}
    qui est indépendant de $ \theta$ donc $S_n$ est bien une statistique exhaustive en $ \theta$. 
\end{example}


\subsection{Minimalité et complétude}

On souhaite définir des statistiques dites minimales, contenant uniquement l'information nécessaire 
sur $ \theta$ à partir de l'échantillon. Pour cela, on pose : 
    $$ \mathcal{A}^* := \bigcap_{T \text{ exhaustives}} \sigma(T) $$ 
appelée la \emph{tribu exhaustive minimale}. Elle contient donc toute l'information apportée par 
$(X_1, \dots, X_n)$ sur $ \theta$. 

\begin{definition}[Statistique exhaustive minimale]
    On dira d'une statistique exhaustive $S^*$ qu'elle est \emph{minimale} si $ \sigma(S^*) = \mathcal{A}^*$. 
\end{definition}

\begin{proposition}[Caractérisation de $ \mathcal{A}$]
    Sous $ H_1$, pour $ \theta_0 \in \Theta$ fixé, on a : 
        $$ 
            \mathcal{A} = \sigma \left\{
                \frac{L_\theta}{L_{\theta_0}} (X_1, \dots, X_n) \mid \theta \in \Theta
            \right\}
        $$
\end{proposition}

\begin{proposition}[Statistiques exhaustives minimales]
    Soit $S$ une statistique exhaustive telle que : 
        \begin{equation}
            S(x_1, \dots, x_n) = S(y_1, \dots, y_n) \quad \iff \frac{L_\theta(x_1, \dots, x_n)}{L_\theta(y_1, \dots, y_n)}
            \text{ indépendants de } \theta
        \end{equation}
    alors $S$ est \emph{minimale}. 
\end{proposition}

\begin{example}[Loi exponentielle]
    Soit la densité de la loi exponentielle $ f_\theta(x) = \theta e^{ - \theta x} $. On cherche à déterminer $S^*$. 
    Posons $ S_n = \sum_{i=1}^{n} X_i$. On a : 
        $$ L_\theta(X_1, \dots, X_n) = \prod_{i=1}^{n} f_\theta(X_i) = \theta^n e^{- \theta S_n} $$ 
    Pour tout $ \theta_0 \not = \theta_1$, on a :
        $$ 
            \frac{L_{\theta_1}}{L_{\theta_2}}(X_1, \dots, X_n) = \exp (( \theta_0 - \theta_1) S_n + n(\log \theta_1 - \log \theta_0))
        $$ 
    Donc $S$ est fonction mesurable du rapport $ L_{\theta_1}/L_{\theta_2}$, elle est donc minimale. 
\end{example}

\begin{definition}[Statistique Complète (ou Totale)]
    Une statistique $T$ est dite \emph{complète} si pour toute fonction mesurable $h$ telle que $\mathbb{E}_\theta[h(T)]$ existe :
    $$ \forall \theta \in \Theta, \quad \mathbb{E}_\theta[h(T)] = 0 \implies h(T) = 0 \quad \mathbb{P}_\theta\text{-p.s.} $$
\end{definition}

\begin{theorem}[Lien entre Complétude et Minimalité]
    Si une statistique $T$ est \emph{exhaustive} et \emph{complète}, alors elle est \emph{exhaustive minimale}.
\end{theorem}

\begin{remark}
    Dans le cas des \emph{modèles exponentiels} (voir plus loin), la statistique naturelle $T(X)$ est toujours exhaustive 
    et complète (si l'espace des paramètres $\Theta$ est d'intérieur non vide). C'est le moyen le plus rapide en exercice 
    pour obtenir une minimale.
\end{remark}

\subsection{Statistiques Libres}

\begin{definition}[Statistique Libre]
    Une statistique $S$ est dite \emph{libre} si sa loi ne dépend pas de $\theta$.
    $$ \forall \theta, \theta' \in \Theta, \quad \mathbb{P}_\theta^S = \mathbb{P}_{\theta'}^S $$
\end{definition}

Cette notion est fondamentale pour isoler le "bruit" du "signal" ($\theta$). Le lien avec l'exhaustivité est donné par le théorème suivant, très utile pour montrer l'indépendance de deux variables ou pour prouver qu'une statistique n'est pas exhaustive.

\begin{theorem}[Théorème de Basu]
    Si $T$ est une statistique \textbf{exhaustive complète} et $S$ est une statistique \textbf{libre}, alors $T$ et $S$ sont \textbf{indépendantes} sous toute loi $\mathbb{P}_\theta$.
\end{theorem}

\begin{example}
    Dans un modèle gaussien $\mathcal{N}(\mu, \sigma^2)$ avec $\sigma^2$ connu, la moyenne empirique $\bar{X}_n$ est exhaustive complète. La variance empirique $S_n^2$ (ou toute valeur centrée $X_i - \bar{X}_n$) est libre. Par le théorème de Basu, $\bar{X}_n$ et $S_n^2$ sont indépendantes.
\end{example}

% ===================================================================================================================
% Information de Küllback-Leibler

\section{Information de Kullback-Leibler}

\begin{definition}[Information de Kullback-Leibler]
    L'information de \emph{Kullback-Leibler} (ou entropie relative) entre deux lois du modèle est définie par : 
    \begin{equation}
        K( \theta_0 \mid \theta_1) = \mathbb{E}_{\theta_0} \left[\log \frac{f_{\theta_0}(X)}{f_{\theta_1}(X)} \right] = 
        \int_\mathcal{X} f_{\theta_0} \log \left( \frac{f_{\theta_0}}{f_{\theta_1}}\right) \; d \mu 
    \end{equation}
\end{definition}

On parle de \emph{divergence} ou de \emph{contraste} plutôt que de distance, car cette application n'est pas 
symétrique ($K(\theta_0|\theta_1) \neq K(\theta_1|\theta_0)$) et ne respecte pas l'inégalité triangulaire.

\begin{proposition}[Propriétés fondamentales]
    Soit $\mathcal{P}_\Theta$ un modèle dominé.
    \begin{itemize}
        \item \emph{Indépendance :} $K( \theta_0 \mid \theta_1)$ ne dépend pas du choix de la mesure dominante $\mu$. 
        \item \emph{Positivité (Inégalité de Gibbs) :} $ K( \theta_0 \mid \theta_1) \geqslant 0$ (découle de l'inégalité de Jensen).
        \item \emph{Séparation :} $ K( \theta_0 \mid \theta_1) = 0 \iff \mathbb{P}_{\theta_0} = \mathbb{P}_{\theta_1}$. 
        \item \emph{Lien Identifiabilité :} Si le modèle est identifiable, $K( \theta_0 \mid \theta_1) = 0 \iff \theta_0 = \theta_1$.
    \end{itemize}
\end{proposition}

% \begin{remark}[Lien avec Fisher]
%     Rappelons que sous les hypothèses de régularité, l'information de Fisher est la courbure locale de la divergence de Kullback-Leibler :
%     $ \mathcal{I}(\theta_0) = \nabla_{\theta}^2 K(\theta_0 | \theta) \big|_{\theta = \theta_0} $.
% \end{remark}


% ===================================================================================================================
% Informations Quantitatives

\section{Informations Quantitatives}

\subsection{Vecteur Score}

Durant toute cette section, nous serons amenés à calculer le gradient de log-vraissemblance. 
Pour simplifier les calculs et étudier leurs propriétés, nous allons les définir sous 
la forme de \emph{vecteurs scores}. 

\begin{definition}[Vecteur Score]
    Soit $(H_{1,2})$, le vecteur aléatoire score d'un échantillon $X$ est défini par le gradient de la 
    log-vraissemblance de l'échantillon. On a donc : 
    \begin{equation}
        \boxed{S_\theta(X) = \nabla_\theta \log f_\theta(X) = \frac{1}{f_\theta(X)} \nabla f_\theta(X). }
    \end{equation}
    C'est une \emph{variable aléatoire}. 
\end{definition}

\begin{prop}[Espérance]
    Sous $(H_{1, 2, 3})$, les vecteurs scores donc centrés. 
        $$ i.e \quad (H_{1, 2, 3}) \Longrightarrow \forall \theta \in \Theta, \; \mathbb{E}_\theta(S_\theta(X)) = 0_p $$
\end{prop}

\begin{proof}
    Supposons $(H_{1, 2, 3})$, on a alors que $ \int_\mathcal{X} f_\theta \; d \mu = 1$ car c'est une densité. 
    Donc par $(H_{1, 2, 3})$, on a : 
        \begin{align*}
            & \int_\mathcal{X} \frac{ \partial}{ \partial \theta_k} f_\theta \; d \mu = 0 \\ 
            \iff & \int_\mathcal{X} \frac{1}{f_\theta(x)} \frac{\partial}{\partial \theta_k} f_\theta(x) \times 
                \underbrace{ f_\theta(x) \; d \mu(x)}_{d \myP_\theta(x)} = 0 \\ 
            \iff & \int_\mathcal{X} \frac{ \partial }{\partial \theta_k} \log f_\theta \; d \myP_\theta = 0 \\ 
            \iff & \mathbb{E}_\theta \left( \frac{\partial}{\partial \theta_k} \log f_\theta \right) = 0 
        \end{align*}
\end{proof}

\subsection{Modèles exponentiels}

Voyons à présent un modèle régulier très utilisé : le \emph{modèle exponentiel}. 

\begin{definition}[Modèle Exponentiel]
    Un modèle $( \mathcal{X}, \mathcal{A}, \mathcal{P}_\Theta)$ sera dit \emph{exponentiel} si : 
        $$ \boxed{\forall \theta \in \Theta, \quad f_\theta(x) = k_\theta h(x) \exp \left( \langle T(x), \gamma(\theta) \rangle\right).} $$
    Avec $ \gamma( \Theta) \subset \R^q$, $T : \mathcal{X} \longrightarrow \R^q$ et 
    $h : \mathcal{X} \longrightarrow \R$ la fonction de changement de mesure. 

    Il sera qualifié \emph{de plein rang} si $p=q$. On dira aussi qu'il est \emph{canonique} sir $ p=q$ et 
    $ \gamma$ est l'identité. On appellera enfin $T(X)$ la \emph{statistique privilégiée du modèle}. 

    Par soucis de clarté, on notera souvent $ \kappa_\theta = \log k_\theta$. 
\end{definition}

\begin{theorem}[Régularité]
    Un modèle exponentiel canonique de plein rang vérifie $(H_{1, 2, 3, 4})$ et $ k_\theta \in \mathcal{C}^\infty$.
\end{theorem}

\begin{proposition}
    Les modèles canoniques sont très pratiques à manipuler. On a :  
        $$ f_\theta(x) = k_\theta h(x) \exp \left( \langle T(x), \theta \rangle\right) $$ 
    Après calcul, on obtient que $S_\theta(X) = \nabla_\theta \log k_\theta + T(X) \in \R^p$. 
    Or, on sait que $ \mathbb{E}_\theta(S_\theta(X)) = 0$, donc : 
        $$ \mathbb{E}_\theta(T(X)) = - \nabla_\theta \kappa_\theta. $$
\end{proposition}

\begin{prop}[Cas général]
    Dans un cadre plus général, un modèle exponentiel de plein rang admet des scores de la forme : 
        $$ S_\theta(X) = J_{ \gamma}^t \times T(X) + \nabla_\theta \kappa_\theta $$ 
    avec $ \displaystyle J_{\gamma} = \left( \left( \frac{\partial \gamma_i}{ \partial \theta_j} (\theta)\right)_{i,j}\right)$ 
    ma matrice Jacobienne $ p \times p$ de $ \gamma : \Theta \longrightarrow \R^p$. 
    Si $J_{\gamma}$ est définie positive, alors : 
        $$ \mathbb{E}_\theta(T(X)) = - \left( J_\gamma ^{-1}\right)^t \times \nabla_\theta \kappa_\theta $$
\end{prop}

\subsection{Information de Fisher}


Passons à la définition principale de cette partie : l'information de Fisher. Elle quantifie 
l'information contenue sur $ \theta$ dans l'acte d'observer $X$. 

\begin{definition}[Information de Fisher]
    Sous $(H_{1,2})$, la \emph{quantité d'information de Fisher} en $ \theta$ apportée par l'acte d'observer 
    une réalisation $X( \omega)$ sous $ \myP_\theta$ est : 
    \begin{equation}
        \mathcal{I}_\theta(X) = \mathbb{E}_\theta(S_\theta(X) \times S_\theta^t(X))
    \end{equation}
    Si $(H_4)$ est vérifiée, alors $ \mathcal{I}_\theta(X) = \V_\theta(S_\theta(X))$, la 
    matrice de covariance du vecteur score. 
\end{definition}

\begin{prop}[Information de Fisher]
    \begin{itemize}
        \item Sous $(H_{1,2})$, $ \mathcal{I}_\theta(X)$ est intrinsèque au modèle (i.e elle ne dépend plus de 
        la mesure dominante $ \mu$). 
        \item Sous $(H_{1,2,3,4})$, on a : 
        $$ 
            \mathcal{I}_\theta(X) = - \mathbb{E}_\theta(J_{S_\theta(X)}) 
            = - \mathbb{E}_\theta \left( \left( \frac{\partial ^2 \log f_\theta(X)}{\partial \theta_i \partial \theta_j}\right)_{i,j}\right)
        $$
        avec $J_{S_\theta(X)}$ la Jacobienne du score. 
    \end{itemize}
\end{prop}

Ces propriétés seront donc privilégiées pour le calcul pratique de l'information de Fisher. 

\begin{prop}[Cas du modèle exponentiel canonique]
    Dans le cas d'un modèle exponentiel canonique, on a $ \mathcal{I}_\theta(X) = \V_\theta(T(X)) = - H_{ \kappa_\theta}$ 
    la hessienne de $ \kappa_\theta = \log k_\theta$. 
    Donc si $ p = 1$, alors $ \mathcal{I}_\theta(X) = - \frac{ \partial ^2}{ \partial \theta^2} \log k_\theta$. 
\end{prop}

\begin{proposition}
    Comme dit précédement, les modèles exponentiels de plein rang sont très utiles. 
    Ils permettent de calculer simplement l'information de Fisher à partir des fonctions 
    définissant la densité. On a donc : 
        \begin{align*}
            \mathcal{I}_\theta(X) &\underset{\text{def}}{=} J_{\gamma}^t \times \V_\theta(T(X)) \times J_\gamma(\theta) \\ 
            &\underset{\text{prop}}{=} - H_{ \kappa_\theta} - H_{\gamma} \left( J_{\gamma}^{-1}\right)^t \nabla_\theta \kappa_\theta
        \end{align*}
    Si $p = 1$, 
        $$ \mathcal{I}_\theta(X) = ( \gamma'(\theta))^2 \times V_\theta(T(X))
            = \frac{- k_\theta'}{k_\theta} \left( \frac{k_\theta ''}{k_\theta'} + \frac{ \gamma_\theta ''}{\gamma_\theta'} 
            - \frac{k_\theta'}{k_\theta}\right)
        $$
\end{proposition}

\begin{example}
    Soit le modèle $X \sim \mathcal{N}( \theta_1, \theta_2^2)$, avec 
    $ 
        \begin{cases}
            \theta_1 = \mathbb{E}_\theta(X) \\
            \theta_2 = \sqrt{\V_\theta(X)}
        \end{cases}
    $. 
    On a donc : 
        $$ f_\theta(x) = \frac{1}{2 \pi} \exp \left( \langle T(x), \gamma(\theta) \rangle + \kappa_\theta\right) $$
    où $ p =q = 2$ et : 
        $$ 
            T(x) = (x, x^2), \quad \kappa_\theta = - \frac{ \theta_1^2}{2 \theta_2^2} - \log \theta_2, \quad 
            \gamma(\theta) = \left( \frac{ \theta_1}{ \theta_2}, - \frac{1}{ 2 \theta_2^2}\right) 
        $$
    On peut donc ensuite appliquer les formules : 
        $$ J_\gamma = 
            \begin{pmatrix}
                \frac{1}{ \theta_2^2} & - \frac{2 \theta_1}{ \theta_2^2} \\
                0 & \frac{1}{ \theta_2^2}
            \end{pmatrix}
            \quad \V_\theta(T(X)) = 
            \begin{pmatrix}
                \theta_2^2 & Cov(X, X^2) \\ 
                Cov(X, X^2) & \V_\theta(X^2)
            \end{pmatrix}
        $$ 
    On a donc : 
        \begin{align*}
            \mathcal{I}_\theta(X) &= Y_\gamma^t \times V_\theta(T(X)) \times J_\gamma \\ 
            &= - \mathbb{E}_\theta \left( \frac{\partial^2}{ \partial \theta_i \partial \theta_j} 
                \left(
                    - \frac{(X- \theta_1)^2}{2 \theta_2^2} - \log \theta_2
                \right)
                \right)_{i,j}
            = 
                \begin{pmatrix}
                    \frac{1}{ \theta_2^2} & 0 \\ 
                    0 & \frac{2}{ \theta_2^2}
                \end{pmatrix}
        \end{align*}
\end{example}


\subsection{Propriétés de l'information de Fisher}

L'information de Fisher possède des propriétés de structure fondamentales qui découlent de la linéarité du score et de la 
géométrie du modèle.

\begin{enumerate}
    \item \textbf{Additivité :} 
    Soient $S$ et $T$ deux statistiques $\mathcal{P}_\Theta$-indépendantes. Si les modèles associés vérifient $(H_{1,2,3})$, l'information jointe est la somme des informations :
    $$ \mathcal{I}_\theta(S, T) = \mathcal{I}_\theta(S) + \mathcal{I}_\theta(T) $$
    \textit{Conséquence :} Dans le modèle d'échantillonnage i.i.d. $\mathcal{P}_\Theta^{\otimes n}$, on a :
    $$ \boxed{\mathcal{I}_\theta(X_1, \dots, X_n) = n \mathcal{I}_\theta(X_1)} $$

    \item \textbf{Décroissance par transfert (Information image) :} 
    Soit $S$ une statistique. Si le modèle image vérifie $(H_{1,2,3})$, le score image est $S_\theta(S) = \mathbb{E}_\theta[S_\theta(X) \mid S]$. Par l'inégalité de Jensen (ou la formule de décomposition de la variance), on a :
    $$ \mathcal{I}_\theta(S) = \mathbb{V}_\theta(\mathbb{E}_\theta[S_\theta(X) \mid S]) \preceq \mathbb{V}_\theta(S_\theta(X)) = \mathcal{I}_\theta(X) $$
    Ainsi, $\mathcal{I}_\theta(X) - \mathcal{I}_\theta(S)$ est une matrice semi-définie positive. On perd de l'information en condensant les données, sauf si la statistique est exhaustive.

    \item \textbf{Conservation par exhaustivité :} 
    Sous $(H_{1,2,3})$, une statistique $S$ est exhaustive si et seulement s'il n'y a pas de perte d'information par transfert :
    $$ \boxed{S \text{ est exhaustive} \iff \mathcal{I}_\theta(S(X_1, \dots, X_n)) = n \mathcal{I}_\theta(X_1)} $$

    \item \textbf{Lien avec le contraste de Kullback :} 
    Sous $(H_{1,2,3,4})$, pour tout $\theta_0 \in \Theta$ fixé, l'information de Fisher est la matrice Hessienne du contraste de Kullback-Leibler (courbure locale de la divergence) :
    $$ \mathcal{I}_{\theta_0}(X) = \text{Hess}_{\theta} \left( K(\theta_0 \mid \theta) \right) \Big|_{\theta=\theta_0} = \left( \left( \frac{\partial^2 K(\theta_0 \mid \theta)}{\partial \theta_i \partial \theta_j} \right)_{i,j} \right)_{\theta=\theta_0} $$
\end{enumerate}